\chapter{Мгновенные платежи в~мире: Pix, UPI, FedNow, SEPA
Instant}\label{ch:instant-payments-world}

СБП из~предыдущей главы~--- одна из~национальных систем мгновенных безотзывных
переводов между банковскими счетами 24/7/365. Такие же системы работают
в~Бразилии (Pix), Индии (UPI), США (FedNow) и~еврозоне (SEPA Instant Credit
Transfer). Для держателя они похожи, но различаются по~четырём признакам: кто
оператор расчёта; задаёт~ли сама схема способ инициации платежа и~формат QR;
какое место у~небанковских участников; есть~ли в~схеме встроенный антифрод. Эти
различия сохраняются в~банковской обработке: сверка, разбор диспутов и~хранение
доказательств у~каждой системы свои.

\section{Pix (Бразилия)}\label{sec:pix-system}

Pix~--- платёжная схема, оператором которой является центральный банк Бразилии
(Banco Central do Brasil, BCB). Pix запущен 16~ноября 2020~года; участие
обязательно для всех финансовых институтов, у~которых более 500~тыс.\
активных счетов клиентов~\cite{pix,pix-participants}.

Расчёт проходит через систему мгновенных платежей SPI (Sistema de Pagamentos
Instantâneos)~--- выделенную RTGS-инфраструктуру под управлением BCB. В~режиме
RTGS (real-time gross settlement, валовой расчёт в~реальном времени) каждый
перевод исполняется индивидуально и~окончательно, без сведения в~нетто-позицию
и~без расчётных окон, в~отличие от~карточного клиринга, где сотни тысяч операций
между парой участников сводятся в~одну нетто-позицию в~конце дня
(см.~раздел~\ref{sec:recon-net-vs-gross}). Каждый участник держит в~SPI
отдельный расчётный счёт; финальность наступает в~момент дебета счёта
плательщика и~кредита счёта получателя~\cite{pix-spi}. Частного клиринга между
участниками нет.

\subsection*{Ключи и~DICT}\label{subsec:pix-keys-dict}

Pix отправляется не~по~номеру счёта напрямую, а~по~ключу: номеру мобильного
телефона, адресу электронной почты, CPF (идентификационный номер физлица), CNPJ
(идентификатор юрлица) или случайному ключу (EVP), который BCB выдаёт как UUID
по~RFC~4122 (в~бразильских банковских интерфейсах его обычно показывают как
строку из~32~шестнадцатеричных символов). Директория DICT (Diretório de
  Identificadores de Contas
Transacionais) под управлением BCB отображает ключ в~реквизиты счёта
и~участника. Перед отправкой банк плательщика обращается к~DICT, получает
реквизиты получателя и~формирует расчётную инструкцию в~SPI.

\subsection*{Инициация: QR-код и~платёжная строка}\label{subsec:pix-init}

BCB определяет форматы инициации Pix в~отдельном
регламенте~\cite{pix-manual-iniciacao}. Статический QR-код многоразовый,
содержит ключ получателя и~опционально сумму. Динамический QR одноразовый,
генерируется под конкретную транзакцию и~содержит ссылку на~платёжную сессию
у~PSP (payment service provider~--- поставщик платёжных услуг) получателя;
кассовое ПО получает обратный вызов после исполнения. Помимо
сканирования, тот же набор полей передаётся как текстовая платёжная строка,
которую копируют в~буфер обмена и~вставляют в~банковское
приложение~\cite{pix-faq}. Платёжная строка нужна, когда PSP-приложение
и~источник кода открыты на~одном устройстве и~сканировать QR камерой
невозможно.

Формат QR-кода Pix задаёт BCB: EMVCo-совместимый BR Code определён на~уровне
самой схемы. Приложения разных банков понимают один и~тот же код, а~подтверждение
плательщика в~разных каналах приводит к~одному и~тому же сообщению в~SPI.

\subsection*{Сроки и~SLA}\label{subsec:pix-timings}

BCB регулирует сроки Pix отдельным регламентом, периодически пересматривая
разрешённые длительности стадий~\cite{pix-manual-tempos}. Документ задаёт
максимальные интервалы на~отдельные этапы~--- от~запроса к~DICT до~завершения
расчёта~--- и~делает SLA (service-level agreement~--- соглашение об~уровне
сервиса) обязательным. Сообщения между участниками и~SPI идут
по~защищённому каналу с~подписями по~сертификатам ICP-Brasil (Infraestrutura de
Chaves Públicas Brasileira~--- национальная инфраструктура открытых ключей);
требования к~подписи и~каналу прописаны в~отдельном регламенте
по~безопасности~\cite{pix-manual-seguranca}.

\subsection*{Возвраты и~MED}\label{subsec:pix-med}

Pix~--- \textit{push}-перевод (платёж, инициируемый и~подтверждаемый со~стороны
  плательщика; противопоставлен \textit{pull}-сценарию карточного списания,
см.~гл.~\ref{ch:open-banking}) с~немедленной финальностью~--- той же
RTGS-финальностью \emph{в~пределах секунд}, что описана в~СБП формулировкой
\enquote{почти сразу} (см.~раздел~\ref{sec:sbp-finality}): разные термины, один
механизм. Поэтому обычный возврат идёт как новая, обратная по~направлению,
операция Pix. Сценарий \enquote{я~отправил не~туда} в~Pix урегулирован отдельно:
получатель обязан рассмотреть запрос на~возврат, но согласие на~него остаётся
за~получателем, если иное не установлено законом или судом.

Для мошеннических операций BCB ввёл специальный механизм возврата MED
(Mecanismo Especial de Devolução): когда у~банка-отправителя есть фактические
основания считать операцию мошеннической, банк через инфраструктуру Pix
инициирует запрос на~блокировку и~возврат суммы у~получателя. Регламент
исполнения Pix описывает MED как отдельную процедуру внутри
схемы~\cite{pix-manual-fluxos}. MED~--- не chargeback и~не наследует его
правовую конструкцию: получателя не~дебетуют принудительно по~решению схемы,
а~проверяют операцию и~могут вернуть средства, если они ещё доступны на~счёте.
Для плательщика итог функционально близок к~карточному chargeback (канал
восстановления средств при~мошенничестве), но процедура, основания
и~распределение рисков отличаются.

MED первой редакции (запуск~--- 2021~год) блокировал и~возвращал средства только
на~первом счёте-получателе; при~типичной схеме обналичивания, когда мошенник
за~секунды распределяет сумму по~цепочке счетов-дропов, на~момент оспаривания
к~первому счёту оставалось около десятой части средств, и~б\'{о}льшая часть
запросов завершалась отказом по~отсутствию остатка. С~02~февраля 2026~года BCB
сделал обязательной редакцию~MED~2.0, регламентированную резолюциями BCB №~493
и~№~546~\cite{bcb-med-2-0}: добровольное подключение открыто с~23~ноября
2025~года, до~10~мая 2026~года действует переходный период без~санкций. MED~2.0
заменяет одноразовую блокировку первого счёта процедурой \enquote{восстановления
средств} (recuperação de valores): директория DICT по~запросу банка-отправителя
выстраивает граф последующих переводов глубиной до~пяти уровней, рассылает
уведомления о~нарушении участникам по~найденному маршруту денег, и~каждый из~них
блокирует остаток на~своём счёте до~разбора. В~MED~2.0 запрос на~оспаривание
подаётся не~через кол-центр, а~через обязательную кнопку внутри банковского
приложения; банк плательщика обязан уведомить банк получателя в~течение
30~минут, целевой срок возврата по~подтверждённой мошеннической цепочке~---
11~дней с~момента оспаривания.

\section{UPI (Индия)}\label{sec:upi-system}

UPI запущен NPCI (National Payments Corporation of India~--- национальная
корпорация платежей Индии) в~апреле 2016~года и~работает поверх IMPS (Immediate
Payment Service)~--- более раннего сервиса 24/7-переводов. Систему регулирует
Резервный банк Индии (RBI) по~закону о~платёжных и~расчётных системах
2007~года~\cite{upi}. UPI~--- набор открытых API для банков и~небанковских
приложений, а~также центральный коммутатор NPCI, который маршрутизирует запросы
и~участвует в~клиринге.

\subsection*{VPA и~четырёхсторонняя модель}\label{subsec:upi-fourparty}

Реквизит платежа в~UPI~--- виртуальный платёжный адрес VPA формата
\texttt{name@psp} (например, \texttt{abhishek@okaxis}). VPA отображается
на~конкретный банковский счёт без раскрытия номера счёта плательщику.
Как DICT в~Pix, VPA избавляет пользователя от~необходимости запоминать
банковские реквизиты.

В~отличие от~карточной четырёхсторонней схемы
(см.~раздел~\ref{sec:ecosystem-four-party}, где сторонами выступают эмитент,
эквайер, держатель карты и~ТСП), в~UPI участвуют четыре других типа сущностей:
банк-плательщик, банк-получатель, PSP-bank (банк-спонсор, который держит
контракт с~NPCI и~отвечает за~TPAP-приложение) и~TPAP (Third-Party Application
Provider~--- поставщик стороннего приложения с~клиентским интерфейсом). PhonePe,
Google~Pay, Paytm~--- крупнейшие TPAP; PSP-bank отвечает перед регулятором
за~работу этих приложений. NPCI выполняет роль центрального коммутатора:
маршрутизирует запросы,
идентифицирует участников через VPA, ведёт клиринг и~расчёт.

\subsection*{Антифрод и~рыночные ограничения}\label{subsec:upi-tpap-cap}

Чтобы не допустить концентрации, NPCI установил лимит: один TPAP не может
превышать 30~\% месячного объёма UPI-транзакций. Процедура мониторинга
закреплена отдельным циркуляром~\cite{npci-upi-tpap-cap}: при~превышении NPCI
обязан вводить ограничительные меры на~новых пользователях TPAP до~возврата доли
под порог. Это рыночное ограничение поверх технического протокола; ничего
подобного в~правилах СБП, Pix, FedNow или SEPA Instant нет.

В~2025~году NPCI ужесточил требования к~использованию
UPI-API~\cite{npci-upi-oc-215,npci-upi-circulars}. Действующий циркуляр
ограничивает частоту вызова вспомогательных API: проверка баланса~--- не более
50~раз в~сутки на~пользователя, запрос списка счетов~--- не более 25, инструкции
AutoPay исполняются тремя обработками в~сутки вместо непрерывного потока.
Обязательны ежегодный независимый аудит и~квартальная самоаттестация участников.
Эти правила установлены не~в~протоколе UPI, а~в~регулировании поверх него; они
сдерживают избыточную нагрузку на~банковские системы со~стороны массовых TPAP
и~смягчают эффект, близкий к~DDoS (distributed denial of service), при~всплесках
обращений.

\subsection*{Сценарии и~отличие от~Pix}\label{subsec:upi-vs-pix}

UPI поддерживает \textit{push} (\enquote{я~плачу}) и~\textit{collect} (запрос
на~оплату: получатель формирует требование, плательщик подтверждает списание).
QR-инфраструктура в~UPI стандартизирована (Bharat~QR~/ UPI~QR), но конкуренция
между приложениями смещает массовый сценарий в~сторону набора VPA
и~\textit{deep-link} (URL, открывающего нужный экран мобильного приложения
с~предзаполненными реквизитами). В~Pix оплата проходит через банковские
приложения и~зависит от~BCB. В~UPI центральный коммутатор принадлежит NPCI,
а~клиентскую часть собирают конкурирующие TPAP-приложения.

\section{FedNow (США)}\label{sec:fednow-system}

FedNow запущен 20~июля 2023~года и~эксплуатируется Федеральными резервными
банками~\cite{fednow-about,fednow-resources}. До~его появления розничный
американский ландшафт делили Fedwire (межбанковский RTGS для крупных сумм,
не~24/7 до~недавнего расширения), ACH (Automated Clearing House~--- пакетные
межбанковские переводы, T+1) и~частный RTP (Real-Time Payments) The Clearing
House (мгновенный, с~2017~года, но не центробанковский). FedNow добавляет
центробанковский сервис, доступный всем держателям мастер-счёта в~ФРС (master
  account~--- расчётный счёт банка непосредственно в~Федеральной резервной системе,
через который проходят расчёты в~деньгах центрального банка).

\subsection*{Правовая основа}\label{subsec:fednow-legal}

FedNow регулируется Subpart~C 12~CFR (Code of Federal Regulations)
Part~210~\cite{reg-j-subpart-c} и~Operating Circular~8~\cite{fednow-oc8}.
Платёжная инструкция в~FedNow определена как сообщение, переданное через сервис;
правила Subpart~C вытесняют непротиворечащие положения Article~4A UCC (Uniform
Commercial Code~--- единообразный торговый кодекс США). FedNow подпадает прямо
под федеральный закон; ответственность участников жёстче, чем у~частного RTP,
который работает по~правилам собственной системы.

Operating Procedures~\cite{fednow-op-procedures} задают требования к~работе:
24-часовой операционный день, 7/365, никаких плановых перерывов. Отдельного
закрытия дня у~FedNow нет: суточная сверка отделена от~расчётного цикла и~проходит
параллельно сервису.

\subsection*{ISO 20022 как родной формат}\label{subsec:fednow-iso20022}

В~отличие от~UPI и~СБП с~их собственными протоколами, FedNow с~момента запуска
построен поверх ISO~20022 (международного стандарта финансовых сообщений,
подробно разобранного в~разделе~\ref{sec:iso20022}). Банк-плательщик передаёт
\texttt{pacs.008} (Payments Clearing and Settlement~--- клиентский кредитовый
перевод от~банка-отправителя к~банку-получателю) в~FedNow; FedNow дебетует счёт
банка-плательщика в~ФРС, кредитует счёт банка-получателя; банк-получатель
зачисляет средства клиенту. Подтверждение возвращается через \texttt{pacs.002}
(статус исполнения соответствующего pacs.008), уведомления клиентам идут через
\texttt{camt.054} (Cash Management~--- уведомление о~зачислении или списании
по~счёту клиента)~\cite{fednow-iso20022}. Префиксы pacs/camt и~соглашения
об~именовании сообщений ISO~20022 разобраны в~разделе~\ref{sec:iso20022}.
Руководство по~внедрению и~песочница для тестирования сообщений размещены
на~платформе документирования сообщений MyStandards (SWIFT~--- Society for
Worldwide Interbank Financial Telecommunication)~\cite{fednow-tech-overview}.

\subsection*{Расчёт и~ликвидность}\label{subsec:fednow-settlement}

Расчёт в~FedNow валовой, в~деньгах центрального банка, через мастер-счёт
участника в~ФРС. Сетевой расчётный риск (риск дефолта расчётчика/контрагента
в~системе нетто-расчёта до~момента финального расчёта) здесь отсутствует:
расчётчиком является ФРС, и~дефолт расчётчика невозможен (в~отличие
от~частного RTP). У~участника возникает другая задача~--- поддерживать
на~мастер-счёте достаточно ликвидности круглосуточно, включая ночи и~выходные.
Параллельно с~запуском FedNow ФРС расширил доступ к~ликвидности (продлённые часы
  дисконтного окна~--- discount window, механизма краткосрочного обеспеченного
  кредитования банков центральным банком~--- и~сервисы внутридневного
кредитования, intraday credit), чтобы участники не~выпадали из~системы в~часы
низкой ликвидности.

\subsection*{Пробелы FedNow: QR, возвраты, сроки}\label{subsec:fednow-gaps}

Своего QR-кода у~FedNow нет: схема не~описывает инициацию~--- ни
о~ссылке для оплаты, ни о~запросе на~оплату с~деталями назначения платежа, ни
о~форматах QR. Инициация оставлена рынку~--- банковским ядрам,
финтех-приложениям, маркетплейсам. Клиентскую часть участники собирают поверх
FedNow самостоятельно или используют надстройку RTP, где запрос на~оплату
и~стандартизированный QR уже определены.

\section{SEPA Instant Credit Transfer (еврозона)}\label{sec:sepa-instant-system}

SCT Inst (SEPA Instant Credit Transfer)~--- мгновенный кредитовый перевод в~евро
внутри SEPA (Single Euro Payments Area, единой зоны платежей в~евро),
определённый правилами Европейского платёжного совета (European Payments
Council, EPC). Целевое сквозное время~--- 10~секунд, 24/7/365. В~еврозоне
работают две расчётные платформы: центробанковская TIPS (TARGET Instant Payment
  Settlement, оператор Eurosystem~--- ECB и~национальные центробанки
еврозоны)~\cite{ecb-tips} и~частная RT1 (Real-Time~1, оператор EBA Clearing~---
инфраструктурный дочерний оператор Euro Banking Association). Участник PSP
подключается к~одной из~них; между ними настроена взаимная достижимость
(reachability~--- возможность достичь любого участника другой платформы
без~отдельного подключения к~ней).

\subsection*{IPR: обязательность через регулирование}\label{subsec:sepa-ipr}

До~2024~года участие в~SCT Inst было добровольным; за~восемь лет после запуска
(ноябрь 2017~года) покрытие так и~не стало полным. Это изменилось с~принятием
Regulation (EU) 2024/886~--- регламента о~мгновенных платежах (Instant Payments
Regulation, IPR)~\cite{eu-2024-886,ecb-ipr}. Регламент опубликован в~Официальном
журнале ЕС 19~марта 2024~года, вступил в~силу 8~апреля 2024~года и~ввёл сроки
обязательной реализации:

\begin{itemize}
  \item 9~января 2025~года~--- PSP в~еврозоне обязаны принимать SCT Inst;
  \item 9~октября 2025~года~--- PSP в~еврозоне обязаны отправлять SCT Inst
    и~предоставлять услугу проверки получателя (Verification of Payee, VOP);
  \item 9~января 2027~года~--- PSP вне еврозоны обязаны принимать;
  \item 9~июля 2027~года~--- PSP вне еврозоны обязаны отправлять;
  \item 9~июня 2028~года~--- финальные обязательства, в~том числе по~тарифной
    нейтральности (комиссия за~SCT Inst не выше, чем за~обычный SCT) и~доступу
    небанковских PSP к~платёжным системам.
\end{itemize}

IPR отменяет разрешённую ранее наценку за~SCT Inst поверх обычного SCT
и~открывает прямой доступ к~платёжным системам для лицензированных небанковских
PSP.

\subsection*{Правила SCT Inst 2025}\label{subsec:sepa-rulebook}

EPC опубликовал свод правил SCT Inst 2025 v1.0 в~ноябре 2024~года; документ
вступил в~силу 5~октября 2025~года вместе с~основной волной обязательств IPR
и~в~тот же момент был заменён на~v1.1~\cite{epc-sct-inst-rulebook-2025}. Свод
согласован с~IPR: выровнены поля, сроки и~обработка отказов. Историческое
ограничение по~сумме на~одну операцию (100\,000~EUR) IPR отменяет~--- лимит
определяется правилами расчётной системы и~PSP.

\subsection*{Проверка получателя (VOP)}\label{subsec:sepa-vop}

Параллельно с~основным сводом EPC ввёл отдельную схему проверки получателя
(VOP). Перед инициированием SCT/SCT Inst PSP плательщика спрашивает PSP
получателя, соответствует ли указанное имя получателя реквизитам IBAN
(International Bank Account Number~--- международный номер банковского счёта).
Ответ~--- \emph{совпадение}, \emph{близкое совпадение} (с~предложенным именем),
\emph{нет совпадения} или \emph{проверка невозможна}~--- передаётся плательщику
до~подтверждения платежа~\cite{epc-vop}. Правила EPC задают два срока:
отвечающий PSP должен сформировать ответ предпочтительно за~одну
секунду и~не позднее трёх, а~сквозной таймаут на~стороне инициирующего PSP~---
пять секунд. Центральный каталог PSP под управлением EPC маршрутизирует
VOP-запрос на~нужную точку приёма.

VOP~--- обязательная по~регулированию мера антифрода: IPR требует от~PSP её
реализации с~9~октября 2025~года. Клиентский интерфейс A2A-платежа
(account-to-account, прямой перевод со~счёта на~счёт;
см.~раздел~\ref{sec:a2a-payments}) в~еврозоне обязан явно показывать результат
VOP перед подтверждением, а~PSP плательщика отвечает за~корректное поведение,
если плательщик настаивает на~отправке при~ответе \enquote{нет совпадения}.

\section{Сравнение систем мгновенных платежей}\label{sec:instant-comparison}

Сводная таблица~\ref{tab:instant-comparison} сопоставляет пять систем~--- СБП
из~главы~\ref{ch:sbp} плюс четыре разобранных~--- по~пяти параметрам: оператор,
расчёт, протокол, антифрод, возврат.

\begin{table}[H]
  \caption{Сравнение СБП, Pix, UPI, FedNow и~SCT
  Inst.}\label{tab:instant-comparison}
  \begingroup
  \scriptsize
  \setlength{\tabcolsep}{3pt}
  \renewcommand{\arraystretch}{1.2}
  \begin{tabularx}{\textwidth}{
      @{}B{0.16\textwidth}
      P{0.155\textwidth}
      P{0.155\textwidth}
      P{0.155\textwidth}
      P{0.155\textwidth}
      P{0.155\textwidth}@{}
    }
    \toprule
    \headrow
    \thd{\thd{}}
    & \thd{СБП} & \thd{Pix} & \thd{UPI} & \thd{FedNow} & \thd{SCT Inst} \\
    \midrule
    Оператор & НСПК (Банк России) & BCB & NPCI (под RBI) &
    Федеральные резервные банки & TIPS (Eurosystem) и~RT1 (EBA Clearing) \\
    Расчёт & ЦБ через ОПКЦ & BCB через SPI & коммутатор NPCI
    и~банки-расчётчики & мастер-счёт в~ФРС & TIPS~--- ЦБ-деньги;
    RT1~--- ЦБ-расчёт по~итогам \\
    Запуск & 2019 & ноябрь~2020 & апрель~2016 & 20~июля 2023 & ноябрь~2017 \\
    Протокол & ОПКЦ OpenAPI & проприетарный с~подписями
    по~сертификатам ICP-Brasil & UPI API NPCI & ISO~20022 (pacs.008)
    & ISO~20022 (pacs.008) \\
    QR в~правилах схемы & да (НСПК) & да (BR Code, BCB) & да (UPI QR,
    NPCI) & нет & нет \\
    Обязательность & добровольная для~банков, обязательная
    для~крупных в~ряде сценариев & обязательная с~500~тыс.\ счетов &
    регулируется RBI и~NPCI & добровольная & обязательная по~IPR с~2025 \\
    Антифрод схемы & приостановка по~369-ФЗ & MED & лимит доли 30~\%,
    лимит частоты API & в~правилах схемы нет & VOP (IBAN~+~имя) \\
    Возврат & новая операция и~процедура у~банка & новая операция,
    MED для~фрода & новая операция, диспут через PSP & новая операция
    & новая операция, плюс отзыв по~правилам схемы \\
    \bottomrule
  \end{tabularx}%
  \endgroup
\end{table}

\section{Сходство и~различия}\label{sec:instant-lessons}

ISO~20022 распространён, но не универсален: у~UPI собственный API NPCI,
у~Pix~--- собственный протокол с~подписями по~сертификатам ICP-Brasil, в~СБП~---
ОПКЦ OpenAPI; FedNow и~SCT Inst построены прямо на~ISO~20022.

\textbf{Кто оператор расчёта}. В~Pix и~СБП оператором является центральный банк
страны с~собственной инфраструктурой расчёта. FedNow эксплуатируется самими
Федеральными резервными банками. SCT Inst работает на~двух платформах
(центробанковской TIPS и~частной RT1), между которыми настроена взаимная
достижимость. В~UPI центральный коммутатор принадлежит NPCI, расчёт между
банками идёт под надзором RBI. Расчёт в~деньгах центрального банка убирает
сетевой расчётный риск, но обязывает регулятора поддерживать круглосуточный
сервис и~отвечать за~любые сбои инфраструктуры.

\textbf{Где задаётся инициация}. Pix и~UPI определяют форматы QR-кода
и~эквивалентов (платёжная строка в~Pix, \textit{deep-link} в~UPI) в~правилах
самой схемы. СБП~--- аналогично, через НСПК. FedNow и~SCT Inst оставляют
инициацию рынку: схема видит только финальную инструкцию pacs.008. Когда
клиентский интерфейс встроен в~правила схемы, у~плательщиков разных PSP
одинаковый сценарий оплаты; когда вынесен наружу, PSP конкурируют интерфейсами,
зато совместимость становится отдельной задачей рынка.

\textbf{Роль небанковских участников}. UPI ввёл отдельный класс TPAP и~ограничил его
30~\% объёма~\cite{npci-upi-tpap-cap}, чтобы клиентские приложения оставались
конкурентными. IPR в~еврозоне открывает прямой доступ к~платёжным системам
лицензированным небанковским PSP. В~Pix, СБП и~FedNow небанковские приложения
работают через банк-спонсор, и~отдельного класса участника схемы у~них нет.

\textbf{Антифрод как часть схемы}. VOP в~SCT Inst работает \emph{до} инициации
(предотвращение), MED в~Pix~--- \emph{после} (восстановление). В~СБП банк обязан
на~два дня приостановить перевод, если получатель числится в~базе Банка России
о~переводах без~согласия. В~UPI ограничения работают через регуляторные
циркуляры~\cite{npci-upi-oc-215}: лимит частоты вызова API, лимит доли объёма
для одного TPAP, обязательные аудиты. FedNow антифрод в~правилах схемы не
описывает: ответственность несут банк-отправитель и~банк-получатель.

Кто работает сразу с~несколькими системами мгновенных платежей, не получает
единого клиентского процесса и~единой банковской обработки: сверка, разбор
диспутов, интерфейс и~хранение доказательств настраиваются под каждую систему
отдельно, потому что различаются сами правовые конструкции возврата, лимитов
и~антифрода.
